\newpage
\section{Acta Número 04}
\label{sec:acta04}

\noindent \textbf{Datos de la Sesión}
\begin{itemize}[noitemsep]
    \item \textbf{Modalidad:} Virtual - Google Meet.
    \item \textbf{Hora de Inicio:} 18:00
    \item \textbf{Hora de Fin:} 20:00
\end{itemize}

\noindent \textbf{Orden del Día}
\begin{itemize}[noitemsep]
    \item \textbf{Asistencia.}
    \item \textbf{Planificación de Sprints del Semestre:} Definición del cronograma general y duración de los ciclos de desarrollo.
    \item \textbf{Definición de Eventos y Fases Scrum:} Fechas y duraciones para el Refinamiento (Grooming), Sprint Planning, Ejecución y Aseguramiento de Calidad (QA).
    \item \textbf{Documentación y Herramientas:} Asignación de tiempos para documentación por sprint y planificación de la entrega final del pliego (física y web).
    \item \textbf{Cierre de Sprint y Entregables:} Cronograma de Finalización, Incremento, Deployment, Sprint Review y Sprint Retrospective.
    \item \textbf{Firma y aprobación de los presentes.}
\end{itemize}

\noindent \textbf{Desarrolladores (Socios Fundadores)}
\begin{enumerate}[noitemsep]
    \item Pardo Pozo Juan Diego (Jefe de Proyecto).
    \item Arnez Jaimes Mauricio.
    \item Quispe Flores Camila Belen.
    \item Ariñez Bautista Jim Gabriel.
    \item Medina Rodríguez Mateo Jhoustin.
    \item Molina Maldonado Tomas Manuel.
    \item Zeballos Crespo Jonathan.
    \item Gutierrez Guzmán Angheline.
\end{enumerate}

\noindent \textbf{Fecha de la sesión:}\par
Lunes 16 de marzo de 2026

\newpage
\subsection{Desarrollo del Acta}

\noindent \textbf{Asistencia}\par
Se procedió a la toma de asistencia a la hora acordada, corroborando la presencia en sesión de los 8 socios fundadores de la grupo-empresa Byte Busters S.R.L. Al contar con la totalidad de los miembros, se estableció el quórum necesario para iniciar la toma de decisiones definitivas.

\vspace*{0.5cm}
\noindent\textbf{Planificación de Sprints del Semestre}\par
Como punto central de la reunión, se elaboró el cronograma maestro para el desarrollo del proyecto durante todo el semestre. Tras evaluar las necesidades de las Historias de Usuario y la carga de trabajo, el equipo determinó y aprobó unánimemente que la duración de cada Sprint será exactamente de 2 semanas.

\vspace*{0.5cm}
\noindent \textbf{Definición de Eventos y Fases Scrum}\par
Se procedió a definir las fechas, días y duraciones precisas para las distintas ceremonias y fases operativas del marco de trabajo:
\begin{itemize}[noitemsep]
    \item \textbf{Refinamiento (Grooming) y Planning:} Se calendarizaron las sesiones para el desglose de Epics, estimación de esfuerzo, Definition of Ready y diseño de Mockups. Asimismo, se fijaron las fechas para la Sprint Planning, incluyendo la selección de PBIs, definición del Sprint Goal y creación del Sprint Backlog.
    \item \textbf{Ejecución y QA:} Se definieron las modalidades y días para el Daily Scrum, junto con los periodos de desarrollo de las HU y el Code Review. Paralelamente, se establecieron las fechas para la fase de Aseguramiento de Calidad (QA), abarcando pruebas unitarias e integración, validación de criterios de aceptación y corrección de bugs.
\end{itemize}

\vspace*{0.5cm}
\noindent \textbf{Documentación y Herramientas}\par
Se acordó que la documentación técnica (Backend, Frontend y Base de Datos) se realizará iterativamente en cada sprint, apegándose a los estándares previamente definidos. Dada la magnitud del pliego de especificaciones (aproximadamente 5 documentos), se decidió reservar un bloque de 2 semanas exclusivas al final del cronograma para su consolidación. Se aprobó la propuesta de generar esta documentación en formato físico y web, definiendo en sesión las herramientas a utilizar para optimizar este proceso.

\vspace*{0.5cm}
\noindent \textbf{Cierre de Sprint y Entregables}\par
Finalmente, se planificaron las fechas de cierre para cada iteración. Se agendó la \textbf{Finalización e Incremento} definiendo los días para cumplir la Definition of Done (DoD), generar el incremento y realizar el Deployment. Posteriormente, se fijaron las sesiones de \textbf{Sprint Review} para efectuar la demostración del producto a la consultora TIS (Ing. Corina Justina Flores Villarroel), capturar feedback y actualizar el Backlog, culminando cada ciclo con la \textbf{Sprint Retrospective} para el análisis del proceso.

\newpage
\textbf{Firmas y aprobación de los integrantes}
\begin{center}
    \noindent
    \setlength{\tabcolsep}{0pt}
    
    \begin{tabular}{ccc}
        \espacioFirma{assets/img/firmas/mauricio.jpg}{Mauricio Jaimes Arnez}{13751221} & 
        \espacioFirma{assets/img/firmas/camila.jpg}{Camila Belen Quispe Flores}{13097244} &
        \espacioFirma{assets/img/firmas/jim.jpg}{Jim Gabriel Ariñez Bautista}{9517642} \\[1.0cm]
        
        \espacioFirma{assets/img/firmas/mateo.jpg}{Mateo Medina Rodríguez}{9453809} &
        \espacioFirma{assets/img/firmas/tomas.jpg}{Tomas Molina Maldonado}{8051701} &
        \espacioFirma{assets/img/firmas/jonathan.jpg}{Jonathan Zeballos Crespo}{13067494} \\[1.0cm]
    \end{tabular}

    \vspace{0.2cm} 
    \begin{tabular}{cc}
        \espacioFirma{assets/img/firmas/diego.jpg}{Juan Diego Pardo Pozo}{12747783} & 
        \espacioFirma{assets/img/firmas/angheline.jpg}{Angheline Gutierrez Guzmán}{13751815}
    \end{tabular}
\end{center}

\vspace{0.5cm}
\noindent \textbf{Fecha de última revisión:} 23 de marzo de 2026